\subsection{Construction à partir de l’état des douze fibres cyclotomiques}
\label{nuclear:fibras-app-witt}
Sur les marques d’APP, la multiplication \(a:r\mapsto4r\) dans
\(\ZZ/9\ZZ\) satisfait \(a^3=1\). Les feuillets vérifient
\[
 S(ax,ay)=aS(x,y),\qquad P(ax,ay)=a^{-1}P(x,y),
\]
puisque \(16\equiv7\equiv4^{-1}\pmod9\). Les orbites unitaires
ordonnées sont \(\mathcal O_0=(1,4,7)\), \(\mathcal O_1=(2,8,5)\).
Dans chacune, le module d’augmentation
\(A(\mathcal O)=\ker(\ZZ[\mathcal O]\xrightarrow{\sum}\ZZ)\)
hérite du produit qui rend les \(e_r\) orthonormaux. Dans la base
\(b_1=e_r-e_{4r}\), \(b_2=e_{4r}-e_{7r}\), ses matrices sont
\[
 G_{A_2}=\begin{pmatrix}2&-1\\-1&2\end{pmatrix},\qquad
 C=\begin{pmatrix}0&-1\\1&-1\end{pmatrix},\qquad
 C^{\mathsf T}G_{A_2}C=G_{A_2}.
\]
La forme réticulaire provient donc des orbites des marques. Pour les paires \((1,2),(2,3),(3,1)\) et les feuillets \(\mu\in\{\Sigma,\Pi\}\), on obtient
\[
 W_{\APP}=\bigoplus_{(i,j)}\bigoplus_\mu\bigoplus_{\varepsilon=0}^1
 A(\mathcal O_\varepsilon),\qquad
 \ell^\Sigma_{ij}(x)=x_i+x_j,\quad
 \ell^\Pi_{ij}(x)=x_ix_j\pmod9.
\]
Chaque terme conserve la paire, le feuillet et l'orbite. L'application d'état est
\[
 \Psi:(\ZZ/9\ZZ)^3\longrightarrow W_{\APP},\qquad
 \Psi(x)_{ij,\mu,\varepsilon}=\beta_\varepsilon(\ell^\mu_{ij}(x)),
 \quad
 \beta_\varepsilon(r)=
 \begin{cases}e_r-e_{4r},&r\in\mathcal O_\varepsilon,\\0,&r\notin\mathcal O_\varepsilon.
 \end{cases}
\]
\begin{proposition}[Équivariance et rang du module d’état]
\label{nuclear:psi-rango}
L’action de \(C\) dans les six fibres additives et de \(C^{-1}\) dans
les six fibres multiplicatives vérifie \(\Psi(ax)=\rho(a)\Psi(x)\).
Les images engendrent \(W_{\APP}\otimes\QQ\), de dimension vingt-quatre.
\end{proposition}
\begin{proof}
Les identités des feuillets et \(\beta(4r)=C\beta(r)\) prouvent
l’équivariance. Pour le rang, on ordonne les coordonnées par les trois
couples indiqués, par \(\Sigma,\Pi\), par \(\mathcal O_0,\mathcal O_1\)
et par \(b_1,b_2\). Dans chaque orbite, \(\beta\) prend successivement
\((1,0),(0,1),(-1,-1)\). La matrice des images de ces états,
lus par lignes, a pour déterminant \(9\) :
\[
\begin{array}{rrrr}
(0,0,1)&(0,0,2)&(0,0,4)&(0,0,5)\\
(0,1,0)&(0,1,1)&(0,1,2)&(0,1,3)\\
(0,1,4)&(0,1,5)&(0,1,6)&(0,2,0)\\
(0,2,2)&(0,2,3)&(0,4,0)&(0,5,0)\\
(1,0,1)&(1,0,2)&(1,0,4)&(1,0,5)\\
(1,1,0)&(1,2,0)&(1,4,0)&(1,5,0).
\end{array}
\]
L'élimination successive, en réduisant chaque ligne par les précédentes et en normalisant son premier coefficient non nul, donne les colonnes pivots numérotées à partir de zéro
\[
\begin{split}
&(8,10,9,11,0,12,14,16,13,15,17,2,\\
&\hspace{2em}18,19,1,3,20,22,21,23,4,6,5,7)
\end{split}
\]
et pivots
\[
(1,1,1,-1,1,1,1,1,1,-1,3,1,1,3,1,-1,1,1,1,-1,1,1,1,-1).
\]
Leur produit, multiplié par le signe de la permutation des colonnes,
est \(9\). Le mineur inversible atteste le rang total.
\end{proof}

\subsection{Alignement dodécaphasique avec les positions de Witt}
\label{nuclear:alineamiento-witt}
L'indice de la paire est \(i\in\ZZ/3\ZZ\), celui du feuillet \(\mu\in\{0,1\}\) et celui de l'orbite \(\varepsilon\in\{0,1\}\). L'origine dodécaphasique étant fixée, \(\lambda(i,\mu,\varepsilon)
=4i+2\mu+\varepsilon\pmod{12}\) est bijective : les quatre valeurs de chaque \(i\) forment un bloc et les trois blocs sont disjoints. Soit \(R=\left(\begin{smallmatrix}0&1\\1&0\end{smallmatrix}\right)\).
La multiplication prouve \(RCR=C^{-1}\) et
\(R^{\mathsf T}G_{A_2}R=G_{A_2}\). Pour les trivialisations
\(\iota_b\) et les repères \(P_b\) transportés conjointement avec
l’incidence de Witt, nous posons
\[
g_b=\iota_b^{-1}P_bCP_b^{-1}\iota_b,\qquad
T_{i\mu\varepsilon}=\iota_{\lambda(i,\mu,\varepsilon)}^{-1}
P_{\lambda(i,\mu,\varepsilon)}R^\mu.
\]
Alors
\[
g_{\lambda(i,\mu,\varepsilon)}T_{i\mu\varepsilon}
=T_{i\mu\varepsilon}C^{(-1)^\mu}.
\]
La somme directe est un isomorphisme \(C_3\)-équivariant entre les douze
fibres d’état et les douze positions réticulaires. Elle conserve couple,
feuillet, orbite et forme sous le transport simultané de l’origine,
de l’ordre et de l’orientation. Ces données demeurent dans le recollement réticulaire
et le passage au voisin construits ci-dessous.
